\documentclass[12pt,letterpaper]{article}
\usepackage[utf8]{inputenc}
\usepackage{amsmath,amssymb,amsthm}
\usepackage{graphicx}
\usepackage{enumitem}
\usepackage{fancyhdr}
\usepackage{hyperref}
\usepackage{tikz}
\usepackage{pgfplots}
\pgfplotsset{compat=1.17}
\usepackage{booktabs}
\usepackage{tcolorbox}
\usepackage{xcolor}
\usepackage{multicol}

% ----------------------------------------------------------------
% Colors
% ----------------------------------------------------------------
\definecolor{examplecolor}{RGB}{220,240,255}
\definecolor{theoremcolor}{RGB}{255,245,220}
\definecolor{definitioncolor}{RGB}{240,255,240}
\definecolor{analogycolor}{RGB}{255,255,220}
\definecolor{skillcolor}{RGB}{255,235,245}
\definecolor{contextcolor}{RGB}{245,245,250}

% ----------------------------------------------------------------
% Custom tcolorboxes - DOMAIN ADAPTIVE
% ----------------------------------------------------------------

% ALWAYS USE: Definition (for core concepts)
\newtcolorbox{definition}{
  colback=definitioncolor,
  colframe=green!50!black,
  title=Definition,
  fonttitle=\bfseries
}

% USE FOR: Math theorems, Scientific laws, Historical facts, Musical theory
\newtcolorbox{theorem}{
  colback=theoremcolor,
  colframe=orange!80!black,
  title=Key Principle, % ADAPT: "Theorem" | "Law" | "Principle" | "Rule"
  fonttitle=\bfseries
}

% ALWAYS USE: Examples
\newtcolorbox{example}{
  colback=examplecolor,
  colframe=blue!70!black,
  title=Example,
  fonttitle=\bfseries
}

% ALWAYS USE: Solutions
\newtcolorbox{solution}{
  colback=white,
  colframe=gray!60!black,
  title=Solution,
  fonttitle=\bfseries
}

% ALWAYS USE: 5th Grade Analogies
\newtcolorbox{analogy}{
  colback=analogycolor,
  colframe=purple!70!black,
  title=5th Grade Analogy,
  fonttitle=\bfseries
}

% USE FOR: Skills, techniques, procedures
\newtcolorbox{skillbox}{
  colback=skillcolor,
  colframe=red!70!black,
  title=Skill, % ADAPT: "Technique" | "Method" | "Procedure"
  fonttitle=\bfseries
}

% USE FOR: Historical context, real-world applications
\newtcolorbox{contextbox}{
  colback=contextcolor,
  colframe=blue!50!black,
  title=Context, % ADAPT: "History" | "Application" | "Real World"
  fonttitle=\bfseries
}

% ----------------------------------------------------------------
% Header / Footer
% ----------------------------------------------------------------
\pagestyle{fancy}
\fancyhf{}
\lhead{Exam Experts}
\rhead{Comprehensive Guide}
\cfoot{\thepage}

% ----------------------------------------------------------------
% Title info (EDIT PER UNIT)
% ----------------------------------------------------------------
\title{
  \textbf{Unit X: [TOPIC NAME]}\\[4pt]
  \Large Complete Comprehensive Guide with 5th Grade Analogies\\[4pt]
  \large Exam Experts
}
\date{}

\begin{document}

\maketitle
\tableofcontents
\newpage

% ================================================================
% SECTION 1: FOUNDATIONS - THE "WHAT"
% ================================================================

\section{Introduction to [TOPIC]}

\subsection{What Is [TOPIC]?}

\begin{definition}
A \textbf{[CORE CONCEPT]} is [clear, concise definition].
% ADAPT: Define the central object/concept/skill being taught
% Examples: "exponential function" | "violin bow technique" |
%           "supply and demand" | "photosynthesis"
\end{definition}

[Explain the big picture in 2--3 paragraphs. Why does this topic exist?
What problems does it solve? Where does it fit in the larger field?]

\begin{analogy}
[5th-grade analogy that captures the essence]
% EXAMPLES:
% Math:      "A function is like a vending machine..."
% Violin:    "Vibrato is like wiggling your finger when you wave hello..."
% Chemistry: "Ionic bonding is like magnets sticking together..."
% History:   "The Industrial Revolution was like when your whole school
%             switches from paper to iPads..."
\end{analogy}

% ----------------------------------------------------------------
% SUBSECTION: DOMAIN-SPECIFIC CLASSIFICATION
% ----------------------------------------------------------------

\subsection{Types/Categories/Styles of [TOPIC]}
% ADAPT THIS SECTION BASED ON DOMAIN:
% Math:     Types of functions, equation forms
% Music:    Playing styles, techniques, genres
% Science:  Categories, classes, types of phenomena
% History:  Eras, movements, schools of thought
% Business: Models, strategies, approaches

\begin{enumerate}[label=\arabic*.]
  \item \textbf{[Type 1]:} [Description and when/where it's used]
  \item \textbf{[Type 2]:} [Description and when/where it's used]
  \item \textbf{[Type 3]:} [Description and when/where it's used]
  % Add more as needed
\end{enumerate}

% ----------------------------------------------------------------
% SUBSECTION: KEY CHARACTERISTICS
% ----------------------------------------------------------------

\subsection{Key Characteristics of [TOPIC]}
% ADAPT: Choose characteristics relevant to your domain
% DELETE sections that don't apply

\begin{enumerate}
  \item \textbf{[Characteristic 1]:} [Description]
  % Math:    Domain, Range, Asymptotes
  % Music:   Pitch range, Timbre, Dynamics
  % Science: Properties, Behavior, Conditions
  % History: Time period, Key figures, Impact
  % Business: Market size, Competition, Resources
  \item \textbf{[Characteristic 2]:} [Description]
  \item \textbf{[Characteristic 3]:} [Description]
  \item \textbf{[Characteristic 4]:} [Description]
  \item \textbf{[Characteristic 5]:} [Description]
\end{enumerate}

\begin{example}
[Concrete example illustrating identification of type, characteristics, or key features]
% Math:      "Identify the domain and range of $f(x) = 2^x$"
% Violin:    "Identify the bow technique in this passage: [musical notation]"
% Chemistry: "Identify the type of bond in NaCl"
% History:   "Identify the era and key characteristics of the Renaissance"
\end{example}

\begin{solution}
[Step-by-step solution for the example above]
\end{solution}

% ================================================================
% SECTION 2: FUNDAMENTALS - THE "HOW"
% ================================================================

\section{Working with [TOPIC]}

\subsection{Basic Operations/Skills}
% ADAPT: What are the foundational skills needed?
% Math:     Evaluation, computation, manipulation
% Music:    Reading notation, finger placement, breath control
% Science:  Measurement, observation, calculation
% History:  Analysis, chronology, source evaluation
% Business: Calculation, modeling, decision-making

[Explain the most basic way to work with or apply this topic]

\begin{example}
[Give 3--4 short tasks that practice the fundamental skill]
% Math:      "Evaluate $f(x)$ at $x = 0$, $2$, $-1$"
% Violin:    "Play these three notes using proper bow hold: [notation]"
% Chemistry: "Balance this equation: $\text{H}_2 + \text{O}_2 \to \text{H}_2\text{O}$"
% History:   "Place these events in chronological order: [events]"
\end{example}

\begin{solution}
[Show computations or steps cleanly, one line per step]
\end{solution}

\begin{analogy}
[Tie the fundamental skill to something familiar]
% "Evaluating a function is like following a recipe..."
% "Proper bow hold is like holding a pencil but more relaxed..."
% "Balancing equations is like making sure both sides of a scale are equal..."
\end{analogy}

% ----------------------------------------------------------------
% SUBSECTION: INTERPRETATION & UNDERSTANDING
% ----------------------------------------------------------------

\subsection{Reading and Interpreting [DOMAIN-SPECIFIC TERM]}
% ADAPT: What needs to be "read" or "interpreted"?
% Math:     Parameters, coefficients, notation
% Music:    Notation, dynamics, tempo markings
% Science:  Data, graphs, experimental results
% History:  Primary sources, timelines, maps
% Business: Financial statements, market data, metrics

[Explain how to extract meaning from symbols, notation, or representations]

\begin{example}
[Problem requiring interpretation of notation or representation]
\end{example}

\begin{solution}
[State each parameter or element and what it means in words]
\end{solution}

% ================================================================
% SECTION 3: VISUALIZATION & REPRESENTATION
% (ADAPT OR REMOVE IF NOT APPLICABLE)
% ================================================================

\section{Visualizing [TOPIC]}
% DELETE THIS SECTION IF:
%   - The topic has no visual/spatial component
%   - There are no graphs, diagrams, or representations
% ADAPT FOR:
%   Math:     Graphing, coordinate systems
%   Music:    Musical notation, score reading
%   Science:  Diagrams, models, schematics
%   History:  Maps, timelines, infographics
%   Business: Charts, organizational diagrams

\subsection{Basic Representation}

[Describe the standard way this topic is visually represented]

% Optional TikZ graphic:
% \begin{center}
% \begin{tikzpicture}
%   \begin{axis}[
%     axis lines = center,
%     xlabel = $x$,
%     ylabel = $y$,
%     xmin=..., xmax=...,
%     ymin=..., ymax=...,
%     grid = major,
%     width=10cm,
%     height=7cm
%   ]
%   \addplot[
%     domain=...:...,
%     samples=100,
%     color=blue,
%     very thick
%   ]{};
%   \addlegendentry{$f(x) = \dots$}
%   \end{axis}
% \end{tikzpicture}
% \end{center}

\subsection{Variations and Transformations}
% ADAPT: How do representations change?
% Math:     Transformations (shifts, stretches, reflections)
% Music:    Transposition, arrangement, interpretation
% Science:  Different scales, perspectives, models
% History:  Different map projections, timeline formats
% Business: Different chart types, data visualizations

[Explain how the representation changes and what those changes mean]

\begin{analogy}
[Use a visual analogy kids understand]
% "Shifting a graph is like moving a picture on graph paper..."
% "Transposing music is like singing the same song but starting on a different note..."
% "Changing map projections is like flattening an orange peel different ways..."
\end{analogy}

\begin{example}
[Give a specific representation and ask students to:]
\begin{itemize}
  \item Describe what they see
  \item Identify key features
  \item Explain what changes were made
\end{itemize}
\end{example}

\begin{solution}
[Describe transformations in order; compute at least one key point; state domain, range, etc.]
\end{solution}

% ================================================================
% SECTION 4: APPLICATIONS & CONTEXT
% ================================================================

\section{Applications and Real-World Context}

\subsection{Why [TOPIC] Matters}

[Explain the practical importance and real-world applications]

\begin{contextbox}
\textbf{Historical Context / Real-World Use:}

[Provide context: Who developed this? When? Why was it needed? How is it used today?]
% Math:      "Logarithms were invented to simplify calculations before calculators..."
% Music:     "Vibrato became standard in Western classical music during the Romantic period..."
% Chemistry: "Understanding chemical bonds led to plastics, medicines, and materials..."
% Business:  "Supply and demand theory explains everything from concert ticket prices to gas prices..."
\end{contextbox}

\subsection{Worked Applications}

\begin{example}
[Detailed application problem -- could be word problem, performance task, analysis, etc.]
% Math:     Population growth, compound interest, physics problems
% Music:    Performance technique in a real piece
% Science:  Laboratory scenario, field observation
% History:  Analyzing a primary source, explaining an event's impact
% Business: Case study, financial analysis, market strategy
\end{example}

\begin{solution}
[Set up the model/approach, substitute values or evidence, interpret the answer in a sentence]
\end{solution}

% ================================================================
% SECTION 5: CONNECTIONS & RELATED CONCEPTS
% (ADAPT OR REMOVE IF NOT APPLICABLE)
% ================================================================

\section{Related Concepts and Connections}
% USE THIS SECTION FOR:
%   Math:     Inverse functions, related operations
%   Music:    Related techniques, complementary skills
%   Science:  Related phenomena, connected principles
%   History:  Related events, parallel movements
%   Business: Related strategies, complementary methods

\begin{definition}
[Define the related concept and explain its connection to the main topic]
% "Logarithms are the inverse of exponentials..."
% "Pizzicato is the complementary technique to bowing..."
% "Oxidation is always paired with reduction..."
% "The Renaissance was a response to Medieval culture..."
\end{definition}

\begin{analogy}
[Use a "connection" analogy]
% "Inverse functions are like addition vs.\ subtraction -- they undo each other..."
% "Pizzicato vs.\ bowing is like typing vs.\ writing -- different ways to do the same job..."
\end{analogy}

\subsection{Key Relationships}

[Describe how the main topic and related concepts interact]

\begin{enumerate}
  \item \textbf{[Relationship 1]:} [Description]
  \item \textbf{[Relationship 2]:} [Description]
  \item \textbf{[Relationship 3]:} [Description]
\end{enumerate}

% ================================================================
% SECTION 6: ADVANCED TECHNIQUES & PROPERTIES
% ================================================================

\section{Important Techniques and Properties}
% ADAPT: What are the "rules of the game"?
%   Math:     Algebraic properties, laws, identities
%   Music:    Performance techniques, practice methods
%   Science:  Laboratory techniques, analytical methods
%   History:  Analysis frameworks, historiographical methods
%   Business: Strategic frameworks, analytical tools

\begin{skillbox}
\textbf{[Technique/Property Name]:}

[Explain the technique, when to use it, and why it works]
\end{skillbox}

\begin{theorem}
\textbf{[Key Principle/Rule]:}

[State the rule formally]
% Math:      "Product of powers: $x^a \cdot x^b = x^{a+b}$"
% Music:     "Proper bow distribution ensures even tone across the stroke"
% Chemistry: "Law of Conservation of Mass: Matter cannot be created or destroyed"
% History:   "Primary sources provide direct evidence; secondary sources provide interpretation"
\end{theorem}

\begin{analogy}
[Explain at least one property in kid-friendly language]
% "Multiplying inside the exponent means adding outside..."
\end{analogy}

\begin{example}
[Have students apply, expand, condense, or manipulate using the properties/techniques]
\end{example}

\begin{solution}
[Show the algebraic or analytical steps, naming which property is used where]
\end{solution}

% ================================================================
% SECTION 7: COMMON MISTAKES & TROUBLESHOOTING
% ================================================================

\section{Common Mistakes and How to Avoid Them}

\subsection{Mistake \#1: [Description]}

\textbf{Why it happens:} [Explanation]

\textbf{How to avoid it:} [Strategy]

\subsection{Mistake \#2: [Description]}

\textbf{Why it happens:} [Explanation]

\textbf{How to avoid it:} [Strategy]

\subsection{Mistake \#3: [Description]}

\textbf{Why it happens:} [Explanation]

\textbf{How to avoid it:} [Strategy]

\subsection{Mistake \#4: [Description]}

\textbf{Why it happens:} [Explanation]

\textbf{How to avoid it:} [Strategy]

% ================================================================
% SECTION 8: PRACTICE PROBLEMS
% ================================================================

\newpage

\section{Practice Problems}
\label{sec:practice}

\subsection*{Instructions}

For each group, the problem numbers match the section numbers of the notes.
Show all work and present your final answer clearly.

% ---------- SECTION 1: FOUNDATIONS ----------

\subsection{Section 1: Foundations}

\textbf{Problem 1.1}

[Problem aligned with definitions and identification]

\vspace{12pt}

\textbf{Problem 1.2}

[Problem aligned with types/categories]

\vspace{12pt}

\textbf{Problem 1.3}

[Problem aligned with characteristics]

\vspace{12pt}

% ---------- SECTION 2: FUNDAMENTALS ----------

\subsection{Section 2: Fundamentals}

\textbf{Problem 2.1}

[Problem aligned with basic operations/skills]

\vspace{12pt}

\textbf{Problem 2.2}

[Problem aligned with interpretation]

\vspace{12pt}

% ---------- SECTION 3: VISUALIZATION ----------

\subsection{Section 3: Visualization}
% SKIP IF SECTION 3 WAS DELETED

\textbf{Problem 3.1}

[Problem aligned with representation]

\vspace{12pt}

\textbf{Problem 3.2}

[Problem aligned with transformations/variations]

\vspace{12pt}

% ---------- SECTION 4: APPLICATIONS ----------

\subsection{Section 4: Applications}

\textbf{Problem 4.1}

[Application/word problem/real-world scenario]

\vspace{12pt}

\textbf{Problem 4.2}

[Application/word problem/real-world scenario]

\vspace{12pt}

% ---------- SECTION 5: CONNECTIONS ----------

\subsection{Section 5: Connections}
% SKIP IF SECTION 5 WAS DELETED

\textbf{Problem 5.1}

[Problem showing relationship between concepts]

\vspace{12pt}

% ---------- SECTION 6: ADVANCED ----------

\subsection{Section 6: Advanced Techniques}

\textbf{Problem 6.1}

[Problem requiring technique/property application]

\vspace{12pt}

\textbf{Problem 6.2}

[Multi-step problem combining multiple techniques]

\vspace{12pt}

% ---------- MIXED REVIEW ----------

\subsection{Mixed Review}

\textbf{Problem MR.1}

[Challenging problem integrating multiple sections]

\vspace{12pt}

\textbf{Problem MR.2}

[Challenging problem integrating multiple sections]

\vspace{12pt}

\textbf{Problem MR.3}

[Challenging problem integrating multiple sections]

% ================================================================
% SECTION 9: SOLUTIONS TO PRACTICE PROBLEMS
% ================================================================

\newpage

\section{Solutions to Practice Problems}

\subsection{Solutions for Section 1: Foundations}

\textbf{Problem 1.1}

\begin{solution}
[Full worked solution]
\end{solution}

\textbf{Problem 1.2}

\begin{solution}
[Full worked solution]
\end{solution}

\textbf{Problem 1.3}

\begin{solution}
[Full worked solution]
\end{solution}

% ----------------------------------------------------------------

\subsection{Solutions for Section 2: Fundamentals}

\textbf{Problem 2.1}

\begin{solution}
[Full worked solution]
\end{solution}

\textbf{Problem 2.2}

\begin{solution}
[Full worked solution]
\end{solution}

% ----------------------------------------------------------------

\subsection{Solutions for Section 3: Visualization}
% SKIP IF SECTION 3 WAS DELETED

\textbf{Problem 3.1}

\begin{solution}
[Full worked solution]
\end{solution}

\textbf{Problem 3.2}

\begin{solution}
[Full worked solution]
\end{solution}

% ----------------------------------------------------------------

\subsection{Solutions for Section 4: Applications}

\textbf{Problem 4.1}

\begin{solution}
[Full worked solution]
\end{solution}

\textbf{Problem 4.2}

\begin{solution}
[Full worked solution]
\end{solution}

% ----------------------------------------------------------------

\subsection{Solutions for Section 5: Connections}
% SKIP IF SECTION 5 WAS DELETED

\textbf{Problem 5.1}

\begin{solution}
[Full worked solution]
\end{solution}

% ----------------------------------------------------------------

\subsection{Solutions for Section 6: Advanced Techniques}

\textbf{Problem 6.1}

\begin{solution}
[Full worked solution]
\end{solution}

\textbf{Problem 6.2}

\begin{solution}
[Full worked solution]
\end{solution}

% ----------------------------------------------------------------

\subsection{Solutions for Mixed Review}

\textbf{Problem MR.1}

\begin{solution}
[Full worked solution]
\end{solution}

\textbf{Problem MR.2}

\begin{solution}
[Full worked solution]
\end{solution}

\textbf{Problem MR.3}

\begin{solution}
[Full worked solution]
\end{solution}

% ================================================================
% APPENDIX: QUICK REFERENCE
% ================================================================

\appendix

\section{Quick Reference Guide}

\subsection{Key Definitions}

\begin{itemize}
  \item \textbf{[Term 1]:} [Brief definition]
  \item \textbf{[Term 2]:} [Brief definition]
  \item \textbf{[Term 3]:} [Brief definition]
  % Add all major terms
\end{itemize}

\subsection{Important Formulas/Rules/Techniques}
% ADAPT: This could be:
%   Math:     Formula sheet
%   Music:    Fingering chart, technique summary
%   Science:  Key equations, procedural checklist
%   History:  Timeline, key figures
%   Business: Framework diagram, calculation formulas

\begin{enumerate}
  \item [Key item 1]
  \item [Key item 2]
  \item [Key item 3]
\end{enumerate}

\subsection{Study Checklist}

Before moving to the next unit, make sure you can:

\begin{itemize}
  \item[$\square$] Define [core concept] in your own words
  \item[$\square$] Identify different types/categories of [topic]
  \item[$\square$] Perform basic operations/skills confidently
  \item[$\square$] Interpret and read [notation/symbols/representations]
  \item[$\square$] Apply [topic] to real-world problems
  \item[$\square$] Recognize common mistakes and avoid them
  \item[$\square$] Connect [topic] to related concepts
  \item[$\square$] Use advanced techniques when appropriate
\end{itemize}

\end{document}
